% !TeX root = closed-systems-from-comparison-completeness.tex
% ============================================================
\chapter*{Notation and Conventions}
\phantomsection
\addcontentsline{toc}{chapter}{Notation and Conventions}
\label{ch:notation-conventions}
% ============================================================
This section records the notation used repeatedly across the monograph.
Local notation introduced for a single construction is defined in the
chapter where it is used.  When a symbol is reused in a later realized
setting, the text states the realization map or quotient through which
the reuse is justified.

\section*{Primitive comparison data}
\begin{description}[leftmargin=3.1cm,style=nextline]
\item[\(U\)]
The underlying nonempty set of primitive states in a comparison world.
No topology, metric, smooth structure, probability measure, or linear
structure is assumed on \(U\) at the primitive level.
\item[\(\mathcal C\)]
The family of intrinsic binary comparison predicates
\(c:U\times U\to\{0,1\}\).
\item[\((U,\mathcal C)\)]
A comparison world.  The primitive definition is given in
\cref{def:intro-comparison-world} and used throughout the foundational
chapters.
\item[\(G=\Aut(U,\mathcal C)\)]
The intrinsic automorphism group of the comparison world.  Its elements
preserve the comparison predicates and determine the diagonal redundancy
used in the quotient semantics.
\item[\(X_A,X_B\)]
The two factors appearing after the product decomposition forced by
comparison completeness in the relevant closed regime.
\item[\(X\)]
A state space or comparison-state carrier when no factor distinction is
being emphasized.
\end{description}

\section*{Quotients and reports}
\begin{description}[leftmargin=3.1cm,style=nextline]
\item[\(\pi\)]
A quotient map, usually from a representative-level object to the
corresponding physical quotient.
\item[\(\Phys\)]
The physical quotient of representative data by the intrinsic diagonal
redundancy.  In the basic two-factor setting,
\[
        \Phys=(X_A\times X_B)/G .
\]
\item[\(R\)]
A coherent report.  Reports are required to factor through the relevant
quotient when they are to count as physical or semantic content.
\item[\(\sim\)]
An equivalence relation.  Its precise meaning is local to the chapter:
orbit equivalence, semantic equivalence, quotient equivalence, or
gauge-equivalence is stated at the point of use.
\item[\(X/G\)]
The orbit quotient of a \(G\)-action on \(X\).  The symbol denotes
quotient semantics, not the assertion that \(X\) was primitive.
\item[\(X\sslash G\)]
A quotient notation used when the quotient is being read with additional
stack, protocol, or gauge-reduction structure.
\end{description}

\section*{Boolean and refinement notation}
\begin{description}[leftmargin=3.1cm,style=nextline]
\item[\(\mathsf B\), \(\mathsf B_A\), \(\mathsf B_B\)]
Boolean algebras of comparison-generated propositions or their
factor-specific versions.
\item[\(\UF(\mathsf B)\)]
The Stone space of ultrafilters on a Boolean algebra \(\mathsf B\).
\item[\(\varprojlim\)]
An inverse limit.  In the refinement tower it denotes the passage from
finite-stage comparison data to the global limit object.
\item[\(\pr\)]
A projection map in a product, inverse system, or quotient diagram.
\item[\(\Idl\)]
An ideal operator, typically in a Boolean or algebraic filtration
setting.
\end{description}

\section*{Groupoid, path, and transport notation}
\begin{description}[leftmargin=3.1cm,style=nextline]
\item[\(\mathcal G(\mathcal I)\)]
The action groupoid or protocol groupoid associated with an index
system \(\mathcal I\).
\item[\(\Ob\), \(\Mor\)]
Object and morphism sets of a category or groupoid.
\item[\(\Path\)]
A path object or path groupoid, depending on the local construction.
\item[\(\Hol\)]
Holonomy or loop transport data.
\item[\(F^k\)]
The descending transport filtration.  The quotient \(F^2/F^3\) is the
first stable visible obstruction layer in the main transport stack.
\item[\(\gr\)]
The associated graded object of a filtration.
\item[\(\ncl{S}\)]
The normal closure generated by a subset \(S\).
\end{description}

\section*{Smooth realization notation}
\begin{description}[leftmargin=3.1cm,style=nextline]
\item[\(M\)]
A smooth manifold appearing only after a realization has been introduced.
It is not primitive comparison data.
\item[\(TM\), \(T^*M\)]
The tangent and cotangent bundles of a realized smooth manifold.
\item[\(\nabla\)]
A connection in a smooth realization.
\item[\(\LC\)]
Levi--Civita transport or the Levi--Civita connection, when the metric
realization has been fixed.
\item[\(\Riem\), \(\Ric\), \(\Scal\)]
The Riemann curvature tensor, Ricci tensor, and scalar curvature of a
smooth realization.
\item[\(\Diff(M)\)]
The diffeomorphism group of \(M\).
\item[\(\Jtwo\)]
The second-jet comparison map used in the intrinsic--smooth interface.
\item[\(\Curv\)]
A curvature operator or curvature class, defined locally in the
realized transport setting.
\end{description}

\section*{Linear, Hilbert, phase, and matter-sector notation}
\begin{description}[leftmargin=3.1cm,style=nextline]
\item[\(\PhysH\)]
The Hilbert-side physical quotient or physical Hilbert sector after the
Hilbert realization has been introduced.
\item[\(\Obs\)]
An observable algebra or observable sector, with the precise algebra
specified in the relevant chapter.
\item[\(\Chan\)]
A channel operator or channel class, especially in scalar-channel and
matter-sector arguments.
\item[\(\Ad\)]
The adjoint action or an adjoint transport operation.
\item[\(\End\)]
An endomorphism algebra.
\item[\(\rank\)]
Rank in the algebraic or vector-bundle setting under discussion.
\item[\(\dimdiamond\)]
Intrinsic causal-diamond growth dimension.
\end{description}

\section*{Logical-status notation}
\begin{description}[leftmargin=3.1cm,style=nextline]
\item[\textup{(SP1)--(SP5)}]
The clauses of Closed-World Admissibility in
\cref{sp:closed-world}.
\item[\textup{(T)}]
The torsion condition in the intrinsic--smooth interface.
\item[\textup{(D)}]
The detectability axiom used to rule out ghost sectors in faithful
smooth realization.
\item[\textup{(U)}]
The scale-coherence condition selecting the maximal readable sector in
the state-side reconstruction extension.
\item[\textup{(R1)--(R3)}]
The additional realization data for the state-side reconstruction
extension.
\end{description}

\section*{Conventions}
\begin{enumerate}[label=\textup{(\roman*)}]
\item A structure called \emph{intrinsic} is recovered from comparison
data and its admissible quotient or transport structure.
\item A structure called \emph{realized} has been introduced through an
explicit smooth, Hilbert, phase, observer-side, or quantitative reading.
\item A statement called \emph{conditional} is proved under its stated
interface or boundary data; the condition is part of the result.
\item A quotient statement is always read modulo the equivalence relation
or symmetry named in the local construction.
\item Chapter references identify where the formal theorem or definition
is proved.  This front matter is a guide to usage, not a replacement for
the chapterwise definitions.
\end{enumerate}
